\title[Israel 2: Éxodo--Josué]{\bfseries{Israel, el pueblo de Dios}}
\subtitle{Estudio 2. Del Éxodo a la muerte de Josué}
\author{}\institute{}\date{}
\titlegraphic{\includegraphics[width=0.85\linewidth,height=0.38\textheight,keepaspectratio]{portada.jpg}\\[-1pt]
  {\fontsize{5}{6}\selectfont\color{sepia}\textit{Nota.} Elaboración propia: Éx 19:4b y 19:6a en texto hebreo consonántico (BHS).}}
\begin{document}
\frame{\titlepage}

% ===================================================
\section{Introducción}
\begin{frame}{¿Qué hace de esclavos un pueblo?}
\begin{columns}[T,onlytextwidth]
\begin{column}{0.48\textwidth}
\begin{cita}Vosotros me seréis un reino de sacerdotes y una nación santa.\fuente{(Éx 19:6)}\end{cita}
\begin{cita}De ellos son la adopción, la gloria, los pactos, la promulgación de la ley, el culto y las promesas.\fuente{(Ro 9:4)}\end{cita}
\begin{clave}Salir de Egipto no basta: en el Sinaí, tres meses después, Dios revela \textbf{para qué} liberó a Israel.\end{clave}
\end{column}
\begin{column}{0.48\textwidth}
\centering
\begin{tikzpicture}[node distance=0.32cm]
\node[bloqueazul,text width=4.6cm] (a) {Redención antes de la ley\\{\tiny\mdseries Éx 20:2}};
\node[bloque,text width=4.6cm,below=of a] (b) {\textbf{Pacto} en el Sinaí: vocación sacerdotal\\{\tiny Éx 19:4--6}};
\node[bloque,text width=4.6cm,below=of b] (c) {\textbf{Presencia}: Dios en medio del pueblo\\{\tiny Éx 25:8}};
\node[bloque,text width=4.6cm,below=of c] (d) {\textbf{Testimonio}: la Torá como sabiduría ante los pueblos\\{\tiny Dt 4:6}};
\node[bloqueoro,text width=4.6cm,below=0.45cm of d] (e) {\textbf{Ruptura y promesa}: la etapa que rompe el pacto guarda los textos de restauración más claros de la Torá\\{\tiny Éx 34:6--7; Dt 4:25--31; 30:1--10}};
\foreach \u/\v in {a/b,b/c,c/d,d/e}{\draw[flecha] (\u) -- (\v);}
\end{tikzpicture}
\elabpropia
\end{column}
\end{columns}
\end{frame}

\begin{frame}{Los ejes de la serie en esta etapa}
\centering
\begin{tikzpicture}
 \foreach \x/\n/\t/\d/\r in {0/1/{Nación santa}/{¿Qué significa la fórmula del Sinaí y cómo se vivió, se quebró y se preservó?}/{Éx 19:4--6; 25:8; Nm 3:12; Dt 4:5--8; Jos 24},
   4.25/2/{Restauración}/{¿Qué hace Dios cuando Israel rompe el pacto recién hecho, y qué anuncia?}/{Éx 32--34; Nm 14; Dt 4:25--31; 30:1--10; 32},
   8.5/3/{Proyección}/{¿Cómo releen los profetas y Pablo el Éxodo, el desierto y Deuteronomio?}/{Is 63; Jer 31:31--34; Os 2:15; Ro 10--11}}{
  \node[caja,anchor=north] (b\n) at (\x,0) {\parbox[t][3.3cm][t]{3.5cm}{\RaggedRight
    {\small\textbf{\textcolor{ladrillo}{\t}}}\\[5pt]\small \d\par\vfill{\scriptsize\color{sepia}(\r)}}};
  \node[circle,fill=cafe,text=piedra,font=\scriptsize\bfseries,inner sep=1.6pt] at (b\n.north west) {\n};}
 \draw[flecha] (b1.east) -- (b2.west);
 \draw[flecha] (b2.east) -- (b3.west);
\end{tikzpicture}
\elabpropia
\vspace{6pt}
\begin{notalam}[Sobre las citas]
Traducción propia desde el texto hebreo (BHS) y griego (NA28); la LXX según Rahlfs-Hanhart. Si la numeración hebrea difiere, va entre corchetes: Dt 29:4 [29:3].
\end{notalam}
\end{frame}

% ===================================================
\section{Cronología}
\begin{frame}{Secuencia de la etapa}
\centering
\begin{tikzpicture}[x=0.78cm]
\draw[line width=1.6pt,cafe,-{Latex[length=3mm]}] (0,0) -- (15.7,0);
\foreach \x/\t/\b in {0.6/{Pascua y paso del mar}/{Éx 12--15},
                      2.7/{Sinaí: pacto y ley}/{Éx 19--24},
                      4.8/{Becerro y tabernáculo}/{Éx 32--40},
                      6.9/{Cades: los espías}/{Nm 13--14},
                      9.0/{Moab: Balaam, Baal-peor}/{Nm 22--25},
                      11.1/{Deuteronomio; muerte de Moisés}/{Dt 1:3; 34},
                      13.2/{Jordán, Jericó, Hai}/{Jos 3--8},
                      15.1/{Siquem; muerte de Josué}/{Jos 24}}{
  \fill[oro] (\x,0) circle (4pt);
  \node[above,font=\scriptsize\bfseries,text=ladrillo,text width=1.6cm,align=center] at (\x,0.2) {\t};
  \node[below,font=\tiny,text=sepia,text width=1.6cm,align=center] at (\x,-0.2) {\b};}
\draw[decorate,decoration={brace,amplitude=5pt,mirror},sepia] (0.4,-0.9) -- (11.3,-0.9) node[midway,below=5pt,font=\scriptsize\itshape]{Cuarenta años: del mar al Jordán (Nm 14:33--34)};
\draw[decorate,decoration={brace,amplitude=5pt,mirror},sepia] (12.9,-0.9) -- (15.3,-0.9) node[midway,below=5pt,font=\scriptsize\itshape]{La tierra};
\end{tikzpicture}
\notafig{Elaboración propia. Secuencia narrativa; las fechas absolutas están en debate.}
\vspace{8pt}
\begin{columns}[T,onlytextwidth]
\begin{column}{0.48\textwidth}
\begin{clave}\footnotesize Caleb tenía 40 años en Cades y 85 al recibir Hebrón (Jos 14:7, 10): unos siete años de campañas.\end{clave}
\end{column}
\begin{column}{0.48\textwidth}
\begin{notalam}[Duración del gobierno de Josué]
28 años según Seder Olam 12; 25 según Josefo (\emph{Ant.} 5.1.29). El texto bíblico no da la cifra.
\end{notalam}
\end{column}
\end{columns}
\end{frame}

\begin{frame}{El calendario interno del relato}
\centering
\begin{tikzpicture}[x=1cm,y=1cm]
\draw[line width=1.4pt,cafe] (0,0) -- (6.4,0);
\draw[line width=1.4pt,cafe,dash pattern=on 4pt off 3pt] (6.4,0) -- (8.6,0);
\draw[line width=1.4pt,cafe,-{Latex[length=2.5mm]}] (8.6,0) -- (12.0,0);
\node[font=\scriptsize\itshape,text=sepia] at (7.5,0.3) {38 años};
\node[font=\tiny,text=sepia] at (7.5,-0.3) {Dt 2:14};
\foreach \x/\t/\b/\up in {0.5/{Salida de Egipto}/{día 15, mes 1\\Nm 33:3}/1,
  2.3/{Llegada al Sinaí}/{«al tercer mes»\\Éx 19:1}/0,
  4.1/{Tabernáculo}/{día 1, mes 1, año 2\\Éx 40:17}/1,
  5.8/{Cades}/{«un año por cada día»\\Nm 14:34}/0,
  9.3/{Deuteronomio}/{mes 11, año 40\\Dt 1:3}/1,
  11.1/{Cruce del Jordán}/{día 10, mes 1\\Jos 4:19}/0}{
  \node[hito] at (\x,0) {};
  \ifnum\up=1
    \node[above=4pt,font=\footnotesize\bfseries,text=ladrillo,align=center] at (\x,0) {\t};
    \node[above=24pt,font=\scriptsize,text=sepia,align=center] at (\x,0) {\b};
  \else
    \node[below=4pt,font=\footnotesize\bfseries,text=ladrillo,align=center] at (\x,0) {\t};
    \node[below=24pt,font=\scriptsize,text=sepia,align=center] at (\x,0) {\b};
  \fi}
\end{tikzpicture}
\elabpropia
\vspace{12pt}
\begin{notalam}[Precisión cronológica]
\footnotesize El relato fecha con exactitud el primer y el último año del desierto. Israel cruza el Jordán el día 10 del primer mes, la fecha en que Éx 12:3 manda separar el cordero; en Guilgal celebra la Pascua (Jos 5:2--12).
\end{notalam}
\end{frame}

\begin{frame}{Fechas en debate}
\centering
\begin{tikzpicture}[x=0.0245cm,y=1.1cm]
% eje: 1600 -> 1150 a.e.c.
\draw[line width=1.2pt,sepia,-{Latex[length=2.5mm]}] (0,0) -- (460,0);
\foreach \a in {1600,1500,1400,1300,1200}{
  \pgfmathsetmacro{\xx}{1600-\a}
  \draw[sepia] (\xx,-0.08) -- (\xx,0.08);
  \node[below,font=\tiny,text=sepia] at (\xx,-0.08) {\a};}
\node[font=\tiny\itshape,text=sepia,anchor=north east] at (460,-0.12) {a.\,e.\,c.};
% temprana
\node[bloqueazul,anchor=west,font=\tiny\bfseries] at (0,1.25) {Datación temprana};
\draw[line width=3pt,ladrillo!70] (154,0.75) -- (194,0.75);
\node[hito] at (154,0.75) {}; \node[hito] at (194,0.75) {};
\node[etiq,above] at (154,0.85) {Éxodo\\c. 1446};
\node[etiq,above] at (200,0.85) {Entrada\\c. 1406};
% tardía
\node[bloqueazul,anchor=west,font=\tiny\bfseries] at (0,-1.35) {Datación tardía};
\fill[celeste!45] (300,-0.95) rectangle (392,-0.6);
\node[etiq] at (346,-0.78) {Entrada: s. \textsc{xiii}};
\node[hito] at (392,-0.78) {};
\node[etiq,right] at (396,-0.78) {Merneptah\\c. 1208};
% Amarna
\fill[oro!25] (200,0.12) rectangle (300,0.32);
\node[font=\tiny,text=tinta] at (250,0.22) {Amarna (s. \textsc{xiv})};
\end{tikzpicture}
\elabpropia
\vspace{4pt}
\begin{columns}[T,onlytextwidth]
\begin{column}{0.48\textwidth}\scriptsize
\textbf{\textcolor{ladrillo}{Temprana}}: 1~R 6:1 (480 años hasta el 4.º año de Salomón) leído de modo literal; Jue 11:26, 300 años en Transjordania (Bimson, 1978).
\end{column}
\begin{column}{0.48\textwidth}\scriptsize
\textbf{\textcolor{ladrillo}{Tardía}}: la ciudad de Ramsés (Éx 1:11) y Ramsés~II; 480 = 12 generaciones de 40 años (Kitchen, 2003).
\end{column}
\end{columns}
\begin{clave}Ninguna fecha está demostrada. La elección modifica la lectura de los hallazgos (Jericó, Hazor), pero deja intactas la secuencia del relato y su teología.\end{clave}
\end{frame}

% ===================================================
\section{Arqueología}
\begin{frame}{Del mar al Jordán}
\begin{columns}[c,onlytextwidth]
\begin{column}{0.55\textwidth}\centering\includegraphics[width=\linewidth,height=0.84\textheight,keepaspectratio]{mapa_exodo.pdf}\par\vspace{1pt}{\fontsize{4.5}{5.5}\selectfont\color{sepia}\textit{Nota.} Elaboración propia (GSHHS/WDB, Basemap); ubicaciones aproximadas.\par}\end{column}
\begin{column}{0.42\textwidth}\scriptsize\RaggedRight
{\scriptsize\RaggedRight
\tikz[baseline=-0.6ex]\draw[ladrillo,line width=1.5pt,line cap=round](0,0)--(0.6,0);~Ruta tradicional: Ramsés $\to$ Sinaí\par
\tikz[baseline=-0.6ex]\draw[cafe,line width=1.5pt,line cap=round](0,0)--(0.6,0);~Sinaí $\to$ Cades (Dt 1:2)\par
\tikz[baseline=-0.6ex]\draw[oro,line width=1.5pt,line cap=round,dash pattern=on 3.4pt off 1.8pt](0,0)--(0.6,0);~Rodeo de Edom $\to$ llanuras de Moab (Nm 20–22)\par}
\medskip

\textbf{Ruta tradicional} (Nm 33): Ramsés, Sucot, el mar, el Sinaí, Cades y el rodeo de Edom hasta las llanuras de Moab.\par\medskip
\textbf{En debate}: el lugar del paso (\heb{יַם־סוּף}, \emph{yam suf}, «mar de los juncos»; LXX «mar Rojo») y el monte Sinaí (Yébel Musa, norte del Sinaí o Madián).\par\medskip
Ein el-Qudeirat (¿Cades?): solo fortalezas del Hierro, aunque un pueblo de tiendas no deja arquitectura.
\end{column}
\end{columns}
\end{frame}

\begin{frame}{Conquista y reparto}
\begin{columns}[c,onlytextwidth]
\begin{column}{0.44\textwidth}\centering\includegraphics[width=\linewidth,height=0.82\textheight,keepaspectratio]{mapa_tribus.pdf}\par\vspace{1pt}{\fontsize{4.5}{5.5}\selectfont\color{sepia}\textit{Nota.} Elaboración propia (GSHHS/WDB, Basemap); ubicaciones aproximadas.\par}\end{column}
\begin{column}{0.53\textwidth}\scriptsize\RaggedRight
{\scriptsize\RaggedRight
\definecolor{mapacarmesi}{HTML}{D62839}
\definecolor{mapanaranja}{HTML}{F28C28}
\definecolor{mapavioleta}{HTML}{7B2CBF}
\definecolor{mapaturquesa}{HTML}{0F8B8D}
\definecolor{mapaazul}{HTML}{1D4E89}
\definecolor{mapaocre}{HTML}{B5651D}
\tikz[baseline=-0.6ex]{\draw[white,line width=3.2pt,line cap=round](0,0)--(0.62,0);\draw[mapacarmesi,line width=1.8pt,line cap=round,-{Latex[length=1.6mm]}](0,0)--(0.62,0);}~Jericó y Hai (Jos 6–8)\par
\tikz[baseline=-0.6ex]{\draw[white,line width=3.2pt,line cap=round](0,0)--(0.62,0);\draw[mapanaranja,line width=1.8pt,line cap=round,dash pattern=on 3.6pt off 2pt,-{Latex[length=1.6mm]}](0,0)--(0.62,0);}~Gabaón $\to$ sur (Jos 10)\par
\tikz[baseline=-0.6ex]{\draw[white,line width=3.2pt,line cap=round](0,0)--(0.62,0);\draw[mapavioleta,line width=1.8pt,line cap=round,dash pattern=on 1.4pt off 1.6pt,-{Latex[length=1.6mm]}](0,0)--(0.62,0);}~Campaña del norte (Jos 11)\par
\smallskip{\color{black!70}\textbf{Itinerario}}\par
\tikz[baseline=-0.7ex]\node[circle,fill=mapacarmesi,draw=white,line width=0.5pt,text=white,font=\fontsize{5}{5}\selectfont\bfseries,inner sep=0.6pt,minimum size=2.6mm]{1};\,Jericó\quad \tikz[baseline=-0.7ex]\node[circle,fill=mapacarmesi,draw=white,line width=0.5pt,text=white,font=\fontsize{5}{5}\selectfont\bfseries,inner sep=0.6pt,minimum size=2.6mm]{2};\,Hai\quad \tikz[baseline=-0.7ex]\node[circle,fill=mapanaranja,draw=white,line width=0.5pt,text=white,font=\fontsize{5}{5}\selectfont\bfseries,inner sep=0.6pt,minimum size=2.6mm]{3};\,Gabaón\quad \tikz[baseline=-0.7ex]\node[circle,fill=mapanaranja,draw=white,line width=0.5pt,text=white,font=\fontsize{5}{5}\selectfont\bfseries,inner sep=0.6pt,minimum size=2.6mm]{4};\,Laquis\quad \tikz[baseline=-0.7ex]\node[circle,fill=mapavioleta,draw=white,line width=0.5pt,text=white,font=\fontsize{5}{5}\selectfont\bfseries,inner sep=0.6pt,minimum size=2.6mm]{5};\,Hazor\par
\makebox[0.62cm][c]{\tikz[baseline=-0.6ex]\node[circle,fill=white,draw=black!80,line width=0.8pt,inner sep=1.2pt]{};}~Tribus: ubicación aproximada (Jos 13–19)\par}
\smallskip

\textbf{Centro} (Jos 6--9): Jericó, Hai y el pacto con Gabaón; Siquem y el Ebal (8:30--35).\par\smallskip
\textbf{Sur} (Jos 10): de Gabaón a Laquis y Hebrón.\par\smallskip
\textbf{Norte} (Jos 11): Hazor, «cabeza de todos aquellos reinos».\par\smallskip
\textbf{Reparto} (Jos 13--19): nueve tribus y media al oeste del Jordán; Rubén, Gad y media Manasés al este. Leví no recibe territorio: «YHWH es su heredad» (Jos 13:33).\par\smallskip
{\color{sepia}\textbf{Precisión:} los límites tribales se representan con rótulos aproximados.}
\end{column}
\end{columns}
\end{frame}

\begin{frame}{Canaán bajo Egipto}
\centering
\begin{tikzpicture}
 \foreach \x/\n/\t/\d/\r in {0/1/{Cartas de Amarna}/{Ciudades-estado vasallas de Egipto. Abdi-Heba de Jerusalén: los \emph{habiru} saquean las tierras del rey.}/{s. \textsc{xiv} a.\,e.\,c.; EA 286--290; Moran, 1992},
   3.2/2/{¿Habiru = hebreos?}/{No sin más: categoría social (desarraigados, mercenarios), de Mesopotamia a Egipto.}/{\heb{עִבְרִי}: relación en debate},
   6.4/3/{Shasu de \emph{Yhw}}/{Listas de Soleb y Amara West: «tierra de los shasu de \emph{Yhw}», región de Seir (cf. Dt 33:2; Jue 5:4).}/{Lectura discutida},
   9.6/4/{Merneptah}/{«Israel está devastado; su descendencia no existe». Israel con determinativo de pueblo.}/{c. 1208 a.\,e.\,c.}}{
  \node[caja,anchor=north] (b\n) at (\x,0) {\parbox[t][3.2cm][t]{2.6cm}{\RaggedRight
    {\small\textbf{\textcolor{ladrillo}{\t}}}\\[4pt]\scriptsize \d\par\vfill{\tiny\color{sepia}(\r)}}};
  \node[circle,fill=cafe,text=piedra,font=\scriptsize\bfseries,inner sep=1.6pt] at (b\n.north west) {\n};}
\end{tikzpicture}
\elabpropia
\vspace{4pt}
\begin{columns}[T,onlytextwidth]
\begin{column}{0.48\textwidth}
\begin{clave}Límite que impone Merneptah: hacia fines del siglo \textsc{xiii}, Israel ya es un pueblo asentado en Canaán.\end{clave}
\end{column}
\begin{column}{0.48\textwidth}
\begin{notalam}[Balaam fuera de la Biblia]
Deir Alla (s. \textsc{viii}): «Escrito de Balaam, hijo de Beor, el que veía a los dioses» (Hoftijzer y van der Kooij, 1976).
\end{notalam}
\end{column}
\end{columns}
\end{frame}

\begin{frame}{La conquista ante la arqueología}
\centering
\begin{tikzpicture}[every node/.style={font=\tiny}]
\foreach \x/\y/\s/\ta/\tb in {
  0/0/{Jericó}/{Garstang; Wood (1990): destrucción c. 1400}/{Kenyon: c. 1550, casi vacía en el s. \textsc{xiii}; radiocarbono: s. \textsc{xvi} (Bruins y van der Plicht, 1995)},
  6.35/0/{Hai (\heb{הָעַי}, «la ruina»)}/{¿Etiología, Betel o Khirbet el-Maqatir?}/{Et-Tell: sin ciudad en el Bronce Tardío (Callaway, 1968)},
  0/-2.15/{Hazor}/{Ben-Tor: destrucción por Israel (Jos 11:10--13)}/{Fuego del s. \textsc{xiii} por revuelta interna (Zuckerman, 2007)},
  6.35/-2.15/{Monte Ebal}/{Zertal: el altar de Jos 8:30--35}/{Kempinski: torre de vigilancia del Hierro~I}}{
  \node[bloqueazul,anchor=north west,text width=5.75cm,font=\scriptsize\bfseries] (t) at (\x,\y) {\s};
  \node[bloque,anchor=north west,text width=2.65cm,minimum height=1.05cm,font=\tiny] (l) at ([yshift=-3pt]t.south west) {\ta};
  \node[bloqueoro,anchor=north east,text width=2.65cm,minimum height=1.05cm,font=\tiny] (r) at ([yshift=-3pt]t.south east) {\tb};
  \draw[<->,sepia,line width=0.6pt] (l.east) -- (r.west);}
\end{tikzpicture}
\notafig{Elaboración propia. Celeste: lectura que apoya el relato; oro: lectura alternativa.}
\vspace{2pt}
\begin{columns}[T,onlytextwidth]
\begin{column}{0.48\textwidth}
\begin{clave}Josué narra la destrucción de solo tres ciudades (Jericó, Hai, Hazor), seguida de un asentamiento largo e incompleto (Jos 13:1; Jue 1).\end{clave}
\end{column}
\begin{column}{0.48\textwidth}
\begin{notalam}[Sierra central y Ebal]
Hierro I: cientos de aldeas nuevas, casas de cuatro espacios, sin huesos de cerdo (Finkelstein, 1988; Faust, 2006). La «tablilla de maldición» del Ebal (2023) es muy cuestionada.
\end{notalam}
\end{column}
\end{columns}
\end{frame}

\begin{frame}{Modelos del origen de Israel}
\centering\scriptsize\renewcommand{\arraystretch}{1.3}
\begin{tabular}{>{\RaggedRight\arraybackslash}p{2.2cm}>{\RaggedRight\arraybackslash}p{2.9cm}>{\RaggedRight\arraybackslash}p{3.3cm}>{\RaggedRight\arraybackslash}p{3.1cm}}
\toprule
{\color{ladrillo}\bfseries Modelo} & {\color{ladrillo}\bfseries Autores} & {\color{ladrillo}\bfseries Propuesta} & {\color{ladrillo}\bfseries Dificultad}\\\midrule
Conquista militar & Albright; Yadin; Wood; Kitchen & Una campaña destruye ciudades cananeas & Fechas de Jericó y Hai\\
Infiltración pacífica & Alt (1925); Noth & Pastores se asientan poco a poco & Ignora la memoria de guerra\\
Revuelta campesina & Mendenhall (1962); Gottwald (1979) & Campesinos rompen con las ciudades en torno a YHWH & Poca evidencia directa\\
Surgimiento local & Finkelstein; Dever; Faust & Población local con aportes externos & Choca con la memoria unánime de Egipto\\
\bottomrule
\end{tabular}
\elabpropia
\vspace{6pt}
\begin{notalam}[El \emph{jérem} y la retórica de conquista]
El \emph{jérem} (\heb{חֵרֶם}) aparece también en la Estela de Mesa; los relatos antiguos de conquista son hiperbólicos (Younger, 1990). Rahab muestra que la pertenencia a Israel no se definía por la etnia.
\end{notalam}
\end{frame}

% ===================================================
\section{El relato}
\begin{frame}{Pascua y paso del mar}
\begin{columns}[T,onlytextwidth]
\begin{column}{0.50\textwidth}
\centering{\large\heb{מִי־כָמֹכָה בָּאֵלִם יְהוָה}}\par\smallskip
\begin{cita}¿Quién como tú entre los dioses, oh YHWH? ¿Quién como tú, magnífico en santidad?\fuente{(Éx 15:11)}\end{cita}
\begin{notalam}[Hebreo]
\heb{פָּסַח} (\emph{pasaj}): el destructor «pasa de largo» (Éx 12:13); \heb{זִכָּרוֹן}, memorial perpetuo. \heb{יְהוָה יִמְלֹךְ לְעֹלָם וָעֶד}: «YHWH reinará por siempre jamás» (15:18), primera confesión del reinado de Dios.
\end{notalam}
\end{column}
\begin{column}{0.46\textwidth}
\centering
\begin{tikzpicture}[node distance=0.3cm]
\node[bloque,text width=4.3cm] (a) {\textbf{La sangre del cordero}\\{\tiny Éx 12:13; una ley para nativo y extranjero, 12:49}};
\node[bloque,text width=4.3cm,below=of a] (b) {\textbf{El mar}: «YHWH peleará por vosotros»\\{\tiny Éx 14:14; el pueblo «creyó», 14:31}};
\node[bloque,text width=4.3cm,below=of b] (c) {\textbf{El cántico}: Dios reina\\{\tiny Éx 15:11, 18}};
\node[bloqueoro,text width=4.3cm,below=of c] (d) {\textbf{Mirada adelante}: la tierra y el santuario\\{\tiny Éx 15:13, 17}};
\node[bloqueazul,text width=4.3cm,below=of d] (e) {«Nuestra Pascua, Cristo»\\{\tiny\mdseries 1~Co 5:7; Ap 15:3}};
\foreach \u/\v in {a/b,b/c,c/d,d/e}{\draw[flecha] (\u) -- (\v);}
\end{tikzpicture}
\elabpropia
\end{column}
\end{columns}
\end{frame}

\begin{frame}{Sobre alas de águila}
\fotolado{panel_ex19.jpg}{Elaboración propia: Éx 19:5--6 en texto hebreo consonántico (BHS).}
\begin{columns}[c,onlytextwidth]
\begin{column}{0.57\textwidth}
\begin{cita}Vosotros visteis lo que hice a los egipcios, y cómo os tomé sobre alas de águila y os he traído a mí. Ahora, pues, si de veras escucháis mi voz y guardáis mi pacto, seréis mi tesoro especial entre todos los pueblos, porque mía es toda la tierra. Y vosotros me seréis un reino de sacerdotes y una nación santa.\fuente{(Éx 19:4--6)}\end{cita}
\centering
\begin{tikzpicture}[node distance=0.32cm]
\node[bloqueazul,text width=1.95cm,minimum height=1.15cm] (a) {Lo que Dios hizo\\{\tiny\mdseries «os he traído \emph{a mí}»}};
\node[bloque,text width=1.95cm,minimum height=1.15cm,right=of a] (b) {\textbf{La condición}\\{\tiny\heb{אִם־שָׁמוֹעַ תִּשְׁמְעוּ}}\\{\tiny «si de veras escucháis»}};
\node[bloqueoro,text width=1.95cm,minimum height=1.15cm,right=of b] (c) {\textbf{La identidad}\\{\tiny tesoro, reino de sacerdotes, nación santa}};
\draw[flecha] (a) -- (b); \draw[flecha] (b) -- (c);
\end{tikzpicture}
\elabpropia
\end{column}
\begin{column}{0.41\textwidth}\end{column}
\end{columns}
\end{frame}

\begin{frame}{Tres palabras del Sinaí}
\centering
\begin{tikzpicture}
 \foreach \x/\n/\t/\h/\d/\r in {0/1/{Segulá}/{סְגֻלָּה}/{Tesoro personal del rey (1~Cr 29:3). «Porque mía es toda la tierra»: elección sin exclusivismo.}/{LXX \grc{λαὸς περιούσιος}},
   4.25/2/{Mamléjet kohaním}/{מַמְלֶכֶת כֹּהֲנִים}/{Israel es ante las naciones lo que el sacerdote es en Israel: acceso, enseñanza, intercesión.}/{LXX \grc{βασίλειον ἱεράτευμα}},
   8.5/3/{Goy qadosh}/{גּוֹי קָדוֹשׁ}/{Una nación entre las naciones, pero apartada. Santidad de pertenencia que genera conducta (Lv 19:2).}/{LXX \grc{ἔθνος ἅγιον}}}{
  \node[caja,anchor=north] (b\n) at (\x,0) {\parbox[t][3.0cm][t]{3.5cm}{\RaggedRight
    {\small\textbf{\textcolor{ladrillo}{\t}}}\hfill\heb{\h}\\[5pt]\scriptsize \d\par\vfill{\scriptsize\color{sepia}\r}}};
  \node[circle,fill=cafe,text=piedra,font=\scriptsize\bfseries,inner sep=1.6pt] at (b\n.north west) {\n};}
\end{tikzpicture}
\elabpropia
\vspace{6pt}
\begin{notalam}[Trasfondo de \emph{segulá}]
Sarna (1991) relaciona \emph{segulá} con el acadio \emph{sikiltu} y con un texto de Ugarit donde un rey vasallo es la \emph{sglt} de su señor.
\end{notalam}
\end{frame}

\begin{frame}{Del Sinaí a 1 Pedro}
\centering
\begin{tikzpicture}[node distance=0.9cm]
\node[caja,text width=3.3cm,align=center] (h) {\textbf{Hebreo (Éx 19:6)}\\[3pt]\heb{מַמְלֶכֶת כֹּהֲנִים}\\\heb{וְגוֹי קָדוֹשׁ}};
\node[caja,text width=3.3cm,align=center,right=of h] (g) {\textbf{LXX}\\[3pt]\grc{βασίλειον ἱεράτευμα}\\\grc{καὶ ἔθνος ἅγιον}};
\node[caja,text width=3.3cm,align=center,right=of g] (p) {\textbf{1~P 2:9}\\[3pt]\grc{βασίλειον ἱεράτευμα,}\\\grc{ἔθνος ἅγιον}\\+ Is 43:20--21};
\draw[-{Latex[length=2.5mm]},oro!90!ladrillo,line width=1.4pt] (h) -- (g);
\draw[-{Latex[length=2.5mm]},oro!90!ladrillo,line width=1.4pt] (g) -- (p);
\end{tikzpicture}
\elabpropia
\vspace{6pt}
\begin{columns}[T,onlytextwidth]
\begin{column}{0.48\textwidth}
\begin{notalam}[Otras lecturas antiguas]
Targum Onqelos: «reyes, sacerdotes»; Ap 1:6 y 5:10: «reino, sacerdotes» (\grc{βασιλείαν, ἱερεῖς}).
\end{notalam}
\end{column}
\begin{column}{0.48\textwidth}
\begin{clave}Pedro escribe a creyentes dispersos, muchos de origen gentil (1~P 1:18; 4:3). Amplía la vocación de Israel a quienes se unen a ella por el Mesías, como las ramas injertadas de Ro 11:17--24; no la traspasa a otro pueblo.\end{clave}
\end{column}
\end{columns}
\end{frame}

\begin{frame}{El Decálogo y el pacto}
\begin{columns}[T,onlytextwidth]
\begin{column}{0.48\textwidth}
\centering{\heb{אָנֹכִי יְהוָה אֱלֹהֶיךָ אֲשֶׁר הוֹצֵאתִיךָ}}\par\smallskip
\begin{cita}Yo soy YHWH tu Dios, que te saqué de la tierra de Egipto, de casa de esclavos.\fuente{(Éx 20:2)}\end{cita}
\begin{notalam}[Pacto y tratados antiguos]
Primero la gracia, después la obediencia, en el mismo orden que el prólogo histórico de los tratados hititas de vasallaje (Mendenhall, 1954). Weinfeld (1972) halló paralelos de Dt 28 con tratados asirios del s. \textsc{vii}; la fecha de redacción sigue discutida.
\end{notalam}
\end{column}
\begin{column}{0.48\textwidth}
\centering
\begin{tikzpicture}[node distance=0.28cm]
\node[bloqueazul,text width=4.3cm] (a) {Gracia: «te saqué de Egipto»\\{\tiny\mdseries Éx 20:2}};
\node[bloque,text width=4.3cm,below=of a] (b) {\textbf{Respuesta}: \heb{נַעֲשֶׂה וְנִשְׁמָע}, «haremos y obedeceremos»\\{\tiny Éx 24:7}};
\node[bloque,text width=4.3cm,below=of b] (c) {\textbf{Ratificación}: \heb{הִנֵּה דַם־הַבְּרִית}, «he aquí la sangre del pacto»\\{\tiny Éx 24:8; Mt 26:28}};
\node[bloque,text width=4.3cm,below=of c] (d) {\textbf{Testimonio}: la ley como sabiduría «ante los ojos de los pueblos»\\{\tiny Dt 4:6--8}};
\node[bloqueoro,text width=4.3cm,below=of d] (e) {\textbf{Fundamento}: «porque YHWH os amó» y por el juramento a los padres\\{\tiny Dt 7:7--8; Ro 11:28}};
\foreach \u/\v in {a/b,b/c,c/d}{\draw[flecha] (\u) -- (\v);}
\draw[flecha,dashed] (e.east) -- ++(0.25,0) |- (a.east);
\end{tikzpicture}
\elabpropia
\end{column}
\end{columns}
\end{frame}

\begin{frame}{Tabernáculo y Shejiná}
\begin{columns}[T,onlytextwidth]
\begin{column}{0.44\textwidth}
\centering{\large\heb{וְעָשׂוּ לִי מִקְדָּשׁ וְשָׁכַנְתִּי בְּתוֹכָם}}\par\smallskip
\begin{cita}Y me harán un santuario, y habitaré en medio de ellos.\fuente{(Éx 25:8)}\end{cita}
\begin{clave}\footnotesize «La gloria de YHWH llenó el tabernáculo» (Éx 40:34). Pablo: «de ellos es la gloria» (Ro 9:4).\end{clave}
\end{column}
\begin{column}{0.53\textwidth}
\centering
\begin{tikzpicture}[node distance=0.35cm]
\node[caja,text width=1.6cm,align=center] (a1) {\heb{מִשְׁכָּן}\\{\tiny morada}};
\node[caja,text width=1.6cm,align=center,right=of a1] (a2) {\grc{σκηνή}\\{\tiny LXX}};
\node[caja,text width=1.6cm,align=center,right=of a2] (a3) {\grc{ἐσκήνωσεν}\\{\tiny Jn 1:14}};
\node[caja,text width=1.6cm,align=center,below=0.7cm of a1] (b1) {\heb{כַּפֹּרֶת}\\{\tiny tapa del arca}};
\node[caja,text width=1.6cm,align=center,right=of b1] (b2) {\grc{ἱλαστήριον}\\{\tiny LXX Éx 25:17}};
\node[caja,text width=1.6cm,align=center,right=of b2] (b3) {\grc{ἱλαστήριον}\\{\tiny Ro 3:25}};
\foreach \u/\v in {a1/a2,a2/a3,b1/b2,b2/b3}{\draw[-{Latex[length=2mm]},oro!90!ladrillo,line width=1.2pt] (\u) -- (\v);}
\end{tikzpicture}
\elabpropia
\vspace{4pt}
\begin{notalam}[Filología]
\emph{Shajan}, «habitar en tienda», da nombre al \emph{mishkán} y al término rabínico \emph{Shejiná}.
\end{notalam}
\end{column}
\end{columns}
\end{frame}

\begin{frame}{Un sacerdocio en círculos}
\begin{columns}[c,onlytextwidth]
\begin{column}{0.44\textwidth}
\centering
\begin{tikzpicture}
\fill[piedra!80!oro,draw=cafe,line width=0.8pt] (0,0) circle (2.45);
\fill[oro!25,draw=cafe,line width=0.8pt] (0,0) circle (1.55);
\fill[oro!55,draw=cafe,line width=0.8pt] (0,0) circle (0.8);
\node[font=\scriptsize\bfseries,text=tinta,align=center] at (0,0) {Aarón\\{\tiny santuario}};
\node[font=\scriptsize\bfseries,text=tinta,align=center] at (0,1.18) {Leví\\[-2pt]{\tiny campamento}};
\node[font=\scriptsize\bfseries,text=tinta,align=center] at (0,2.0) {Israel\\[-2pt]{\tiny entre las naciones}};
\end{tikzpicture}
\elabpropia
\end{column}
\begin{column}{0.53\textwidth}
\begin{cita}He tomado a los levitas [\ldots] en lugar de todos los primogénitos, porque mío es todo primogénito.\fuente{(Nm 3:12--13)}\end{cita}
\scriptsize
\begin{itemize}
\item Israel, primogénito entre las naciones (Éx 4:22); sus primogénitos, librados en la Pascua y representados por los levitas.
\item Bendición sacerdotal: «pondrán mi nombre sobre los hijos de Israel» (Nm 6:24--27).
\end{itemize}
\begin{notalam}[Arqueología: Ketef Hinom]
Fines s. \textsc{vii}--inicios s. \textsc{vi} a.\,e.\,c.: dos rollitos de plata con esta bendición, los textos bíblicos más antiguos conservados (Barkay et al., 2004).
\end{notalam}
\end{column}
\end{columns}
\end{frame}

\begin{frame}{El becerro de oro}
\begin{columns}[T,onlytextwidth]
\begin{column}{0.52\textwidth}
\centering
\begin{tikzpicture}[node distance=0.26cm]
\node[bloqueazul,text width=5.2cm] (a) {Cuarenta días después de «haremos y obedeceremos», el becerro\\{\tiny\mdseries Éx 32:1--6}};
\node[bloque,text width=5.2cm,below=of a] (b) {Dios propone empezar de nuevo con Moisés {\tiny(32:10)}};
\node[etiq,font=\scriptsize\itshape,text=sepia,below=0.25cm of b] (i) {Intercesión de Moisés (32:11--13)};
\node[bloqueoro,text width=1.5cm,minimum height=1.0cm,below=0.3cm of i] (c2) {Tu nombre ante Egipto};
\node[bloqueoro,text width=1.5cm,minimum height=1.0cm,left=0.18cm of c2] (c1) {Es tu pueblo};
\node[bloqueoro,text width=1.5cm,minimum height=1.0cm,right=0.18cm of c2] (c3) {Juramento a los padres};
\node[bloque,text width=5.2cm,below=0.3cm of c2] (d) {«Perdona [\ldots]; y si no, ráeme de tu libro» {\tiny(32:32)}};
\draw[flecha] (a) -- (b); \draw[flecha] (b) -- (i);
\foreach \c in {c1,c2,c3}{\draw[flecha] (i.south) -- (\c.north); \draw[flecha] (\c.south) -- (d.north);}
\end{tikzpicture}
\elabpropia
\end{column}
\begin{column}{0.44\textwidth}
\begin{cita}Desearía yo mismo ser anatema, separado del Mesías, por amor a mis hermanos.\fuente{(Ro 9:3)}\end{cita}
\centering{\footnotesize\grc{ηὐχόμην γὰρ ἀνάθεμα εἶναι αὐτὸς ἐγώ}}\par\medskip
\begin{clave}\footnotesize Romanos 9--11 se abre con una intercesión de tono mosaico (Moo, 1996).\end{clave}
\end{column}
\end{columns}
\end{frame}

\begin{frame}{El nombre misericordioso}
\fotolado{panel_ex34.jpg}{Elaboración propia: Éx 34:6--7 en texto hebreo consonántico (BHS).}
\begin{columns}[c,onlytextwidth]
\begin{column}{0.57\textwidth}
\begin{cita}¡YHWH, YHWH, Dios compasivo y clemente, lento para la ira y grande en misericordia y fidelidad; que guarda misericordia a millares, que perdona la iniquidad, la rebelión y el pecado!\fuente{(Éx 34:6--7)}\end{cita}
\begin{clave}Se revela \textbf{después} de que Israel rompe el pacto; Dios renueva el pacto y escribe de nuevo las tablas (Éx 34:1, 10).\end{clave}
\begin{notalam}[Hebreo y ecos]
\heb{חֶסֶד} (\emph{jésed}) y \heb{אֱמֶת} (\emph{emet}); \heb{רַחוּם}, de \heb{רֶחֶם}, «seno materno». Citado en Nm 14:18; Sal 86:15; 103:8; Jl 2:13; Jon 4:2; Neh 9:17.
\end{notalam}
\end{column}
\begin{column}{0.41\textwidth}\end{column}
\end{columns}
\end{frame}

\begin{frame}{Cades: perdón y remanente}
\begin{columns}[T,onlytextwidth]
\begin{column}{0.48\textwidth}\footnotesize\RaggedRight
\begin{itemize}
\item El pueblo propone volver a Egipto, es decir, deshacer el Éxodo (Nm 14:4).
\item Moisés cita Éx 34:6--7 (Nm 14:17--19).
\item \heb{סָלַחְתִּי כִּדְבָרֶךָ}: «lo he perdonado conforme a tu palabra» (14:20).
\item Pero la generación incrédula muere en el desierto; sus hijos entran (14:22--31).
\end{itemize}
\end{column}
\begin{column}{0.48\textwidth}
\centering
\begin{tikzpicture}[node distance=0.4cm]
\node[bloque,text width=4.2cm] (a) {\textbf{\textcolor{ladrillo}{Perdón del pueblo}}\\{\tiny Nm 14:20}};
\node[bloque,text width=4.2cm,below=of a] (b) {\textbf{\textcolor{ladrillo}{Juicio de una generación}}\\{\tiny Nm 14:29--34}};
\node[bloqueoro,text width=4.2cm,below=of b] (c) {\textbf{\textcolor{ladrillo}{Remanente: Josué y Caleb}}\\{\tiny Nm 14:24, 30}};
\node[bloqueazul,text width=4.2cm,below=of c] (d) {«Un remanente escogido por gracia»\\{\tiny\mdseries Ro 11:5}};
\draw[flecha] (a) -- (b); \draw[flecha] (b) -- (c); \draw[flecha,dashed] (c) -- (d);
\end{tikzpicture}
\elabpropia
\end{column}
\end{columns}
\end{frame}

\begin{frame}{Balaam y Baal-peor}
\begin{columns}[T,onlytextwidth]
\begin{column}{0.50\textwidth}
\centering{\heb{הֶן־עָם לְבָדָד יִשְׁכֹּן}}\par\smallskip
\begin{cita}He aquí un pueblo que habitará solo, y no será contado entre las naciones.\fuente{(Nm 23:9)}\end{cita}
\centering{\heb{דָּרַךְ כּוֹכָב מִיַּעֲקֹב}}\par\smallskip
\begin{cita}Saldrá una estrella de Jacob, y se levantará un cetro de Israel.\fuente{(Nm 24:17)}\end{cita}
\end{column}
\begin{column}{0.46\textwidth}\scriptsize\RaggedRight
\begin{notalam}[La estrella en la tradición]
Targum Onqelos: «se levantará un rey de Jacob y será ungido el Mesías de Israel». Qumrán: CD \textsc{vii} 18--21; 4Q175. Rabí Aquiba y Bar Kojba, «hijo de la estrella» (TJ Taanit 4:5).
\end{notalam}
\begin{itemize}
\item Baal-peor (Nm 25), por consejo de Balaam (Nm 31:16; Ap 2:14); Finees y el «pacto de sacerdocio perpetuo».
\end{itemize}
\begin{clave}La idolatría no entra recién con Balaam: el episodio pesa (Os 9:10), pero el becerro (Éx 32) es anterior, y Jos 24:14 y Ez 20:7--8 la remontan a Egipto.\end{clave}
\end{column}
\end{columns}
\end{frame}

\begin{frame}{Deuteronomio en Moab}
\begin{columns}[T,onlytextwidth]
\begin{column}{0.48\textwidth}
\centering{\large\heb{שְׁמַע יִשְׂרָאֵל יְהוָה אֱלֹהֵינוּ יְהוָה אֶחָד}}\par\smallskip
\begin{cita}Escucha, Israel: YHWH nuestro Dios, YHWH uno es. Y amarás a YHWH tu Dios con todo tu corazón.\fuente{(Dt 6:4--5)}\end{cita}
\begin{notalam}[Tradición judía]
El Shemá es la confesión que Israel recita dos veces al día.
\end{notalam}
\end{column}
\begin{column}{0.48\textwidth}
\centering
\begin{tikzpicture}[node distance=0.28cm]
\node[bloqueazul,text width=4.3cm] (a) {La Torá predicada a la generación nacida en el desierto};
\node[bloque,text width=4.3cm,below=of a] (b) {\textbf{Renovación del pacto}: bendiciones y maldiciones\\{\tiny Dt 27--30}};
\node[bloque,text width=4.3cm,below=of b] (c) {\textbf{Testigo}: el Cántico de Moisés\\{\tiny Dt 31:19; 32}};
\node[bloqueoro,text width=4.3cm,below=of c] (d) {«Profeta como yo te levantará YHWH»\\{\tiny Dt 18:15; Hch 3:22}};
\node[bloque,text width=4.3cm,below=of d] (e) {«Nunca más se levantó profeta como Moisés»\\{\tiny Dt 34:10}};
\foreach \u/\v in {a/b,b/c}{\draw[flecha] (\u) -- (\v);}
\draw[flecha,dashed] (d) -- (e);
\end{tikzpicture}
\elabpropia
\end{column}
\end{columns}
\end{frame}

\begin{frame}{Josué: la entrada en la tierra}
\centering
\begin{tikzpicture}
 \foreach \x/\n/\t/\d/\r in {0/1/{Jordán y Guilgal}/{El cruce repite el mar; doce piedras «para que todos los pueblos conozcan»; circuncisión y Pascua.}/{Jos 3--5; 4:24},
   3.2/2/{Rahab}/{Una cananea confiesa: «YHWH es Dios arriba en los cielos y abajo en la tierra»; Mt 1:5.}/{Jos 2:11; 6:25},
   6.4/3/{Hai y Acor}/{Acán toma lo consagrado; castigo en el \emph{émeq Ajor}, «valle de la desgracia».}/{Jos 7:24--26},
   9.6/4/{Ebal y Gabaón}/{La Torá leída ante «los extranjeros»; los gabaonitas incorporados por un pacto jurado.}/{Jos 8:30--35; 9:27}}{
  \node[caja,anchor=north] (b\n) at (\x,0) {\parbox[t][3.1cm][t]{2.6cm}{\RaggedRight
    {\small\textbf{\textcolor{ladrillo}{\t}}}\\[4pt]\scriptsize \d\par\vfill{\tiny\color{sepia}(\r)}}};
  \node[circle,fill=cafe,text=piedra,font=\scriptsize\bfseries,inner sep=1.6pt] at (b\n.north west) {\n};}
\end{tikzpicture}
\elabpropia
\vspace{6pt}
\begin{notalam}[El nombre de Josué]
\heb{יְהוֹשֻׁעַ} (\emph{Yehoshúa}, «YHWH salva»), en griego \grc{Ἰησοῦς}: «si Josué les hubiera dado el reposo\ldots» (Heb 4:8).
\end{notalam}
\end{frame}

\begin{frame}{La tierra dada}
\begin{columns}[T,onlytextwidth]
\begin{column}{0.48\textwidth}
\begin{cita}No faltó palabra de todas las buenas promesas que YHWH había hecho a la casa de Israel; todo se cumplió.\fuente{(Jos 21:45)}\end{cita}
\begin{notalam}[Crítica textual]
El Códice de Leningrado omite Jos 21:36--37; algunas ediciones hebreas numeran 21:41--43.
\end{notalam}
\end{column}
\begin{column}{0.48\textwidth}
\centering
\begin{tikzpicture}[node distance=0.3cm]
\node[bloqueazul,text width=4.3cm] (a) {\heb{נַחֲלָה} (\emph{najalá}): heredad};
\node[bloque,text width=4.3cm,below=of a] (b) {No se vende a perpetuidad: «porque la tierra mía es»\\{\tiny Lv 25:23}};
\node[bloque,text width=4.3cm,below=of b] (c) {Repartida por tribus; Leví: «YHWH es su heredad»\\{\tiny Jos 13--21; 13:33}};
\node[bloqueoro,text width=4.3cm,below=of c] (d) {Aun así, «queda mucha tierra por poseer»\\{\tiny Jos 13:1}};
\foreach \u/\v in {a/b,b/c,c/d}{\draw[flecha] (\u) -- (\v);}
\end{tikzpicture}
\elabpropia
\end{column}
\end{columns}
\end{frame}

\begin{frame}{Siquem: el pacto renovado}
\begin{columns}[T,onlytextwidth]
\begin{column}{0.40\textwidth}
\centering{\heb{וְאָנֹכִי וּבֵיתִי נַעֲבֹד אֶת־יְהוָה}}\par\smallskip
\begin{cita}Pero yo y mi casa serviremos a YHWH.\fuente{(Jos 24:15)}\end{cita}
\begin{clave}\footnotesize Siquem cierra el arco: donde Abraham recibió la promesa, Israel renueva el pacto.\end{clave}
\end{column}
\begin{column}{0.56\textwidth}
\centering
\begin{tikzpicture}[x=1cm,y=1cm]
\node[bloque,text width=2.5cm] (g1) at (0,1.5) {\textbf{Siquem}: la promesa a Abraham\\{\tiny Gn 12:6--7}};
\node[bloque,text width=2.5cm] (g2) at (0,0) {\textbf{José} pide que lleven sus huesos a la tierra\\{\tiny Gn 50:25}};
\node[bloqueoro,text width=2.5cm] (j1) at (3.6,1.5) {\textbf{Siquem}: renovación del pacto\\{\tiny Jos 24}};
\node[bloqueoro,text width=2.5cm] (j2) at (3.6,0) {Huesos de José sepultados en Siquem\\{\tiny Jos 24:32}};
\draw[flecha] (g1) -- (j1); \draw[flecha] (g2) -- (j2);
\node[font=\tiny\bfseries,text=ladrillo] at (0,2.35) {Génesis};
\node[font=\tiny\bfseries,text=ladrillo] at (3.6,2.35) {Josué};
\draw[sepia,dashed] (1.8,2.4) -- (1.8,-0.75);
\end{tikzpicture}
\elabpropia
\end{column}
\end{columns}
\end{frame}

% ===================================================
\section{Vocación}
\begin{frame}{Vivida, quebrada, preservada}
\centering\scriptsize\renewcommand{\arraystretch}{1.3}
\begin{tabular}{>{\RaggedRight\arraybackslash}p{3.7cm}>{\RaggedRight\arraybackslash}p{3.3cm}>{\RaggedRight\arraybackslash}p{4.2cm}}
\toprule
{\color{ladrillo}\bfseries Se vivió} & {\color{ladrillo}\bfseries Se quebró} & {\color{ladrillo}\bfseries Se preservó}\\\midrule
«Haremos y obedeceremos» (Éx 24:7) & Becerro de oro (Éx 32) & Intercesión; nombre misericordioso (Éx 34:6--7)\\
La gloria en el tabernáculo (Éx 40:34) & Rebelión de Cades (Nm 14) & Perdón; remanente: Josué y Caleb\\
Torá ante extranjeros; Rahab; Gabaón (Jos 8:35) & Baal-peor (Nm 25) & Finees: «sacerdocio perpetuo» (Nm 25:13)\\
Levitas fieles (Éx 32:26--29) & Acán (Jos 7) & Pacto renovado en Ebal y Siquem\\
\bottomrule
\end{tabular}
\elabpropia
\vspace{6pt}
\begin{tikzpicture}[node distance=0.35cm]
\node[bloque,text width=3.4cm,minimum height=1.1cm] (a) {\textbf{Hacia las naciones} desde el comienzo\\{\tiny Éx 12:49; Dt 4:6}};
\node[bloque,text width=3.4cm,minimum height=1.1cm,right=of a] (b) {\textbf{Santidad por presencia}: «¿en qué se conocerá\ldots sino en que tú andes con nosotros?»\\{\tiny Éx 33:16}};
\node[bloqueoro,text width=3.4cm,minimum height=1.1cm,right=of b] (c) {\textbf{Mediador y remanente} cuando el pueblo falla};
\end{tikzpicture}
\end{frame}

% ===================================================
\section{Tradición judía}
\begin{frame}{Shavuot y la Torá para todos}
\begin{columns}[T,onlytextwidth]
\begin{column}{0.46\textwidth}\footnotesize
\textbf{\textcolor{ladrillo}{Shavuot}}\par\smallskip
Fiesta de primicias (Lv 23:15--21); unida al Sinaí por Éx 19:1. Liturgia: \heb{זְמַן מַתַּן תּוֹרָתֵנוּ}, «tiempo de la entrega de nuestra Torá».\par\smallskip
\begin{notalam}[Talmud]
Los Diez Mandamientos, el 6 de Siván según los sabios, el 7 según rabí Yosé (Shabat 86b). «Cada palabra se dividía en setenta lenguas» (Shabat 88b); posible eco en Hch 2, aunque no demostrado.
\end{notalam}
\end{column}
\begin{column}{0.50\textwidth}
\centering
\begin{tikzpicture}
\node[bloqueazul,text width=1.8cm] (t) at (0,0) {La Torá\\{\tiny\mdseries Dt 33:2}};
\node[bloque,text width=1.5cm] (e) at (-2.1,1.25) {Esaú};
\node[bloque,text width=1.5cm] (m) at (2.1,1.25) {Amón y Moab};
\node[bloque,text width=1.5cm] (i) at (-2.1,-1.25) {Ismael};
\node[bloqueoro,text width=1.5cm] (s) at (2.1,-1.25) {Israel};
\foreach \n in {e,m,i}{\draw[flecha,dashed] (t) -- (\n);}
\draw[flecha] (t) -- (s);
\node[etiq,above=1pt of e] {la rehúsa};
\node[etiq,above=1pt of m] {la rehúsan};
\node[etiq,below=1pt of i] {la rehúsa};
\node[etiq,below=1pt of s] {«haremos y obedeceremos»};
\end{tikzpicture}
\notafig{Elaboración propia a partir de Sifré Deuteronomio 343.}
\vspace{2pt}
\begin{clave}Se dio en el desierto, lugar sin dueño, para que ninguna nación dijera que no tenía parte en ella (Mekhilta, Bajódesh 5). Cf. Éx 19:5: «mía es toda la tierra».\end{clave}
\end{column}
\end{columns}
\end{frame}

\begin{frame}{El calendario de la memoria}
\centering
\begin{tikzpicture}[x=1cm,y=1cm]
\draw[line width=1.4pt,cafe,-{Latex[length=2.5mm]}] (0,0) -- (12.4,0);
\foreach \x/\m in {0.9/Nisán,3.0/Siván,5.2/Tamuz,7.3/Av,10.4/Tishré}{
  \node[font=\tiny\bfseries,text=sepia] at (\x,-0.28) {\m};}
\foreach \x/\d/\t/\c/\up in {0.9/{15 Nisán}/{Pésaj: salida de Egipto}/bloque/1,
  3.0/{6 Siván}/{Shavuot: entrega de la Torá}/bloque/0,
  5.2/{17 Tamuz}/{Se rompen las tablas}/bloqueoro/1,
  7.3/{9 Av}/{Decreto contra la generación de Cades; caen ambos templos}/bloqueoro/0,
  9.6/{10 Tishré}/{Segundas tablas; Yom Kipur}/bloque/1,
  11.4/{Sukkot}/{Cabañas del desierto (Lv 23:43)}/bloque/0}{
  \node[hito] at (\x,0) {};
  \ifnum\up=1
    \node[\c,text width=1.85cm,font=\tiny,anchor=south] at (\x,0.45) {\textbf{\d}\\\t};
    \draw[sepia,line width=0.4pt] (\x,0.08) -- (\x,0.45);
  \else
    \node[\c,text width=1.85cm,font=\tiny,anchor=north] at (\x,-0.5) {\textbf{\d}\\\t};
    \draw[sepia,line width=0.4pt] (\x,-0.08) -- (\x,-0.5);
  \fi}
\end{tikzpicture}
\notafig{Elaboración propia a partir de Mishná, Taanit 4:6; Shabat 86b; Seder Olam 6.}
\vspace{4pt}
\begin{columns}[T,onlytextwidth]
\begin{column}{0.48\textwidth}
\begin{notalam}[9 de Av]
«Llorasteis por nada; os estableceré un llanto por generaciones» (Taanit 29a). Las \emph{trece midot} de Éx 34:6--7 son el «orden de la plegaria» para el perdón (Rosh Hashaná 17b).
\end{notalam}
\end{column}
\begin{column}{0.48\textwidth}
\begin{notalam}[Pésaj y Sukkot en la liturgia]
\emph{Az yashir} y \emph{Mi jamoja} cada mañana; \emph{Dayenu} en el Séder. Cabañas: nubes de gloria (R. Eliézer) o cabañas reales (R. Aquiba) (Sukká 11b); Zac 14:16.
\end{notalam}
\end{column}
\end{columns}
\end{frame}

\begin{frame}{La cadena de la tradición}
\centering{\large\heb{משה קבל תורה מסיני ומסרה ליהושע}}\par\smallskip
\begin{cita}Moisés recibió la Torá desde el Sinaí y la transmitió a Josué; Josué, a los ancianos; los ancianos, a los profetas.\fuente{(Mishná, Avot 1:1)}\end{cita}
\begin{tikzpicture}[node distance=0.55cm]
\node[bloqueazul,minimum width=2cm] (a) {Sinaí};
\node[bloque,minimum width=2cm,right=of a] (b) {Moisés};
\node[bloque,minimum width=2cm,right=of b] (c) {Josué};
\node[bloque,minimum width=2cm,right=of c] (d) {Ancianos};
\node[bloqueoro,minimum width=2cm,right=of d] (e) {Profetas};
\foreach \u/\v in {a/b,b/c,c/d,d/e}{\draw[flecha] (\u) -- (\v);}
\end{tikzpicture}
\elabpropia
\vspace{2pt}
\begin{columns}[T,onlytextwidth]
\begin{column}{0.48\textwidth}\scriptsize
\begin{itemize}
\item Josué escribió su libro y los ocho versículos finales de la Torá (Bava Batra 15a).
\item Rahab, esposa de Josué y antepasada de profetas, entre ellos Jeremías (Meguilá 14b): haggadá.
\end{itemize}
\end{column}
\begin{column}{0.48\textwidth}\scriptsize
\begin{itemize}
\item Las manos de Moisés: «mientras Israel miraba hacia arriba [\ldots] prevalecía» (Mishná, Rosh Hashaná 3:8).
\item Josefo: tabernáculo como imagen del cosmos (\emph{Ant.} 3.7.7); muerte de Moisés (\emph{Ant.} 4.8.48).
\end{itemize}
\end{column}
\end{columns}
\end{frame}

% ===================================================
\section{Restauración}
\begin{frame}{La promesa en la Torá}
\centering\scriptsize\renewcommand{\arraystretch}{1.3}
\begin{tabular}{>{\RaggedRight\arraybackslash}p{2.2cm}>{\RaggedRight\arraybackslash}p{4.7cm}>{\RaggedRight\arraybackslash}p{4.5cm}}
\toprule
{\color{ladrillo}\bfseries Texto} & {\color{ladrillo}\bfseries Promesa} & {\color{ladrillo}\bfseries Proyección}\\\midrule
Éx 32:13; 34:6--10 & Perdón por el juramento a los padres y el nombre misericordioso & Jl 2:13; Jon 4:2; Sal 103:8\\
Nm 14:20--24 & Perdón, juicio de una generación, remanente & Is 10:20--22; Ro 11:5\\
Nm 24:17 & Estrella de Jacob & 4Q175; Mt 2:2\\
Dt 4:25--31 & Destierro y regreso «en los postreros días» & Os 3:5; Jer 29:12--14\\
Dt 30:1--10 & Regreso; corazón circuncidado & Jer 31:33 [31:32]; Ez 36:26; Ro 10:6--8\\
Dt 32:21, 43 & Celos; las naciones con su pueblo & Ro 10:19; 11:11; 15:10\\
\bottomrule
\end{tabular}
\elabpropia
\vspace{6pt}
\begin{clave}En esta etapa la promesa es a la vez respuesta inmediata a la infidelidad (Éx 34; Nm 14) y profecía de un futuro lejano (Dt 4; 30; 32).\end{clave}
\end{frame}

\begin{frame}{Destierro y regreso anunciados}
\begin{columns}[T,onlytextwidth]
\begin{column}{0.46\textwidth}
\centering{\heb{כִּי אֵל רַחוּם יְהוָה אֱלֹהֶיךָ}}\par\smallskip
\begin{cita}En los postreros días volverás a YHWH tu Dios [\ldots]. Porque Dios compasivo es YHWH tu Dios; no te dejará [\ldots] ni se olvidará del pacto que juró a tus padres.\fuente{(Dt 4:30--31)}\end{cita}
\begin{notalam}[Vocabulario profético]
De aquí tomarán los profetas \heb{בְּאַחֲרִית הַיָּמִים} (\emph{be-ajarit ha-yamim}, «en los postreros días»), el verbo \emph{shuv} («volverás») y el \emph{rajum} de Éx 34:6.
\end{notalam}
\end{column}
\begin{column}{0.50\textwidth}
\centering
\begin{tikzpicture}[node distance=0.28cm]
\node[bloqueazul,text width=4.6cm] (a) {\textbf{Promesa}: juramento a los padres\\{\tiny\mdseries Dt 4:31; 7:8}};
\node[bloque,text width=4.6cm,below=of a] (b) {\textbf{Condición}: «amarás a YHWH tu Dios con todo tu corazón»\\{\tiny Dt 6:5; bendiciones y maldiciones, Dt 27--30}};
\node[bloqueoro,text width=4.6cm,below=of b] (c) {\textbf{Ruptura anunciada}: «YHWH os esparcirá entre los pueblos»\\{\tiny Dt 4:27}};
\node[bloque,text width=4.6cm,below=of c] (d) {\textbf{Restauración}: «volverás» y Dios «hará volver tu cautividad»\\{\tiny Dt 4:30; 30:2--3}};
\node[bloqueazul,text width=4.6cm,below=of d] (e) {Corazón circuncidado\\{\tiny\mdseries Dt 30:6}};
\foreach \u/\v in {a/b,b/c,c/d,d/e}{\draw[flecha] (\u) -- (\v);}
\end{tikzpicture}
\elabpropia
\end{column}
\end{columns}
\end{frame}

\begin{frame}{El corazón circuncidado}
\fotolado{panel_dt30.jpg}{Elaboración propia: Dt 30:6 en texto hebreo consonántico (BHS).}
\begin{columns}[c,onlytextwidth]
\begin{column}{0.57\textwidth}
\begin{cita}Y circuncidará YHWH tu Dios tu corazón y el corazón de tu descendencia, para que ames a YHWH tu Dios con todo tu corazón y con toda tu alma, a fin de que vivas.\fuente{(Dt 30:6)}\end{cita}
\centering
\begin{tikzpicture}[node distance=0.3cm]
\node[bloque,text width=1.95cm,minimum height=1.0cm] (a) {\textbf{Mandato}\\{\tiny Dt 10:16}};
\node[bloqueoro,text width=1.95cm,minimum height=1.0cm,right=of a] (b) {\textbf{Promesa}\\{\tiny Dt 30:6}};
\node[bloqueazul,text width=1.95cm,minimum height=1.0cm,right=of b] (c) {Nuevo pacto; corazón nuevo\\{\tiny\mdseries Jer 31; Ez 36:26--27}};
\draw[flecha] (a) -- (b); \draw[flecha] (b) -- (c);
\end{tikzpicture}
\elabpropia
\vspace{2pt}
\scriptsize\RaggedRight
\begin{itemize}
\item Israel «volverá» (\emph{shuv}); Dios «hará volver tu cautividad» (\heb{אֶת־שְׁבוּתְךָ}, 30:3) y lo recogerá «del extremo de los cielos» (30:4).
\item ONE FOR ISRAEL (2025) incluye Dt 4:30--31 entre los textos del regreso nacional esperado en Ro 11:26.
\end{itemize}
\end{column}
\begin{column}{0.41\textwidth}\end{column}
\end{columns}
\end{frame}

\begin{frame}{El Cántico de Moisés}
\begin{columns}[T,onlytextwidth]
\begin{column}{0.46\textwidth}
\centering
\begin{tikzpicture}[node distance=0.22cm]
\node[bloque,text width=4.2cm] (a) {\textbf{Elección}\\{\tiny Dt 32:8--9}};
\node[bloque,text width=4.2cm,below=of a] (b) {\textbf{Cuidado} «como el águila»\\{\tiny 32:11}};
\node[bloqueoro,text width=4.2cm,below=of b] (c) {\textbf{Infidelidad y juicio}; los celos\\{\tiny 32:21}};
\node[bloqueazul,text width=4.2cm,below=of c] (d) {\textbf{Compasión}: «por amor de sus siervos se arrepentirá»\\{\tiny\mdseries 32:36}};
\node[bloque,text width=4.2cm,below=of d] (e) {\textbf{Las naciones} se alegran con su pueblo\\{\tiny 32:43}};
\foreach \u/\v in {a/b,b/c,c/d,d/e}{\draw[flecha] (\u) -- (\v);}
\end{tikzpicture}
\elabpropia
\end{column}
\begin{column}{0.50\textwidth}
\centering\scriptsize\renewcommand{\arraystretch}{1.25}
\begin{tabular}{>{\RaggedRight\arraybackslash}p{1.15cm}>{\RaggedRight\arraybackslash}p{4.3cm}}
\toprule
\multicolumn{2}{l}{\color{ladrillo}\bfseries Dt 32:21}\\\midrule
TM & \heb{וַאֲנִי אַקְנִיאֵם בְּלֹא־עָם}\\
LXX & \grc{κἀγὼ παραζηλώσω αὐτοὺς ἐπ᾽ οὐκ ἔθνει}\\
Ro 10:19 & \grc{ἐγὼ παραζηλώσω ὑμᾶς ἐπ᾽ οὐκ ἔθνει}\\\midrule
\multicolumn{2}{l}{\color{ladrillo}\bfseries Dt 32:43}\\\midrule
TM & \heb{הַרְנִינוּ גוֹיִם עַמּוֹ}\\
LXX = Ro 15:10 & \grc{εὐφράνθητε, ἔθνη, μετὰ τοῦ λαοῦ αὐτοῦ}\\
\bottomrule
\end{tabular}
\elabpropia
\vspace{2pt}
\begin{notalam}[Crítica textual]
Dt 32:8: TM «hijos de Israel»; 4QDeut\textsuperscript{j} «hijos de Dios»; LXX «ángeles de Dios». Dt 32:43: 4QDeut\textsuperscript{q} y LXX más largos (Tov, 2012).
\end{notalam}
\end{column}
\end{columns}
\end{frame}

\begin{frame}{Isaías: el nuevo éxodo}
\begin{columns}[T,onlytextwidth]
\begin{column}{0.48\textwidth}
\centering{\heb{אַיֵּה הַשָּׂם בְּקִרְבּוֹ אֶת־רוּחַ קָדְשׁוֹ}}\par\smallskip
\begin{cita}¿Dónde el que los hizo subir del mar con el pastor de su rebaño? ¿Dónde el que puso en medio de él su Santo Espíritu? [\ldots] El Espíritu de YHWH los hizo descansar.\fuente{(Is 63:11, 14)}\end{cita}
\begin{notalam}[El Santo Espíritu]
Junto con Is 63:10 y Sal 51:11 [51:13], uno de los tres pasajes de la Biblia hebrea que hablan del «Santo Espíritu» de Dios.
\end{notalam}
\end{column}
\begin{column}{0.48\textwidth}
\centering
\begin{tikzpicture}
\node[bloqueazul,text width=1.6cm] (c) at (0,0) {Nuevo éxodo};
\node[bloque,text width=2.0cm] (a) at (-1.55,1.75) {\textbf{Is 63:7--14}\\{\tiny el Éxodo como obra del Espíritu}};
\node[bloque,text width=2.0cm] (b) at (1.55,1.75) {\textbf{Is 43:16--21}\\{\tiny «este pueblo he formado para mí»}};
\node[bloque,text width=2.0cm] (d) at (-1.55,-1.75) {\textbf{Is 11:16}\\{\tiny camino «como el día que subió de Egipto»}};
\node[bloqueoro,text width=2.0cm] (e) at (1.55,-1.75) {\textbf{Is 61:6}\\{\tiny «sacerdotes de YHWH»: Éx 19:6}};
\foreach \n in {a,b,d,e}{\draw[flecha] (c) -- (\n);}
\end{tikzpicture}
\elabpropia
\vspace{2pt}
{\scriptsize \heb{עַם־זוּ יָצַרְתִּי לִי} (Is 43:21; cf. 1~P 2:9).}
\end{column}
\end{columns}
\end{frame}

\begin{frame}{Jeremías: el nuevo pacto}
\begin{columns}[T,onlytextwidth]
\begin{column}{0.48\textwidth}
\begin{cita}[fontupper=\scriptsize\itshape]Haré un nuevo pacto con la casa de Israel y con la casa de Judá. No como el pacto que hice con sus padres el día que los tomé de la mano para sacarlos de la tierra de Egipto [\ldots]. Pondré mi ley en su interior y la escribiré en su corazón.\fuente{(Jer 31:31--33 [31:30--32])}\end{cita}
\begin{notalam}[Destinatarios]
Casa de Israel (reino del norte, ya deportado por Asiria cuando escribe Jeremías) y casa de Judá (reino del sur); véanse los estudios 4 y 5.
\end{notalam}
\end{column}
\begin{column}{0.48\textwidth}
\centering
\begin{tikzpicture}
\node[bloqueazul,text width=1.9cm] (s) at (-1.2,0) {Sinaí};
\node[bloqueazul,text width=1.9cm] (n) at (1.2,0) {Nuevo pacto};
\node[bloque,text width=1.9cm,minimum height=0.8cm] at (-1.2,-0.95) {Tablas de piedra\\{\tiny Éx 34:1}};
\node[bloqueoro,text width=1.9cm,minimum height=0.8cm] at (1.2,-0.95) {El corazón\\{\tiny Jer 31:33; Dt 30:6}};
\node[bloque,text width=4.3cm] at (0,-1.95) {\textbf{Igual}: «mi ley» y los mismos destinatarios};
\node[bloque,text width=4.3cm] at (0,-2.75) {\textbf{Perdón}: el de Éx 34:6--7; Ro 11:27};
\draw[flecha] (s) -- (n);
\end{tikzpicture}
\elabpropia
\vspace{2pt}
\scriptsize
\begin{itemize}
\item \textbf{Jer 2:2--3}: el desierto como noviazgo; «santo era Israel para YHWH, primicias de sus frutos» (cf. Ro 11:16).
\item \textbf{Jer 7:22--23}: en el Éxodo Dios habló primero de «escuchad mi voz» (Éx 19:5).
\end{itemize}
\end{column}
\end{columns}
\end{frame}

\begin{frame}{Oseas: puerta de esperanza}
\centering{\large\heb{וְאֶת־עֵמֶק עָכוֹר לְפֶתַח תִּקְוָה}}\par\smallskip
\begin{cita}La llevaré al desierto y le hablaré al corazón [\ldots], y el valle de Acor por puerta de esperanza; y allí responderá como en el día en que subió de la tierra de Egipto.\fuente{(Os 2:14--15 [2:16--17])}\end{cita}
\vspace{4pt}
\begin{tikzpicture}[node distance=0.9cm]
\node[bloqueoro,text width=3.2cm,minimum height=1.1cm] (a) {\textbf{Jos 7:26}\\[2pt]\mdseries Acán: «valle de la desgracia»};
\node[bloque,text width=3.2cm,minimum height=1.1cm,right=of a] (b) {\textbf{Os 2:15 [2:17]}\\[2pt]«puerta de esperanza»};
\node[bloqueazul,text width=3.2cm,minimum height=1.1cm,right=of b] (c) {Os 1:10 [2:1]; Ro 9:26\\[2pt]\mdseries «hijos del Dios viviente»};
\draw[flecha] (a) -- (b);
\draw[flecha] (b) -- (c);
\end{tikzpicture}
\elabpropia
\vspace{4pt}
\begin{clave}Por donde Israel cayó al comenzar la conquista volverá a entrar (cf. Is 65:10).\end{clave}
\end{frame}

\begin{frame}{Pablo lee esta etapa}
\fotolado{panel_1p29.jpg}{Elaboración propia: 1~P 2:9 en escritura uncial (texto NA28).}
\begin{columns}[c,onlytextwidth]
\begin{column}{0.57\textwidth}\scriptsize
\renewcommand{\arraystretch}{1.18}
\begin{tabular}{>{\RaggedRight\arraybackslash}p{1.45cm}>{\RaggedRight\arraybackslash}p{5.3cm}}
\toprule
{\color{ladrillo}\bfseries Pablo} & {\color{ladrillo}\bfseries Texto de la Torá y lectura}\\\midrule
Ro 10:6--8 & Dt 30:12--14: «cerca de ti está la palabra», aplicada a la palabra de fe\\
Ro 11:8 & Dt 29:4 [29:3]: ojos que no ven «hasta el día de hoy»\\
Ro 10:19; 11:11 & Dt 32:21: los celos para atraer de nuevo a Israel\\
Ro 15:10 & Dt 32:43 LXX: «alegraos, naciones, con su pueblo»\\
1~Co 10:1--11 & A gentiles: «nuestros padres pasaron el mar» (Hays, 1989)\\
\bottomrule
\end{tabular}
\elabpropia
\vspace{2pt}
\begin{notalam}[El arco de Dt 29--30]
Ceguera presente, corazón circuncidado prometido. Dt 30:12--14: cf. Baruc 3:29--30; Targum Neofiti.
\end{notalam}
\end{column}
\begin{column}{0.41\textwidth}\end{column}
\end{columns}
\end{frame}

% ===================================================
\section{Síntesis}
\begin{frame}{Síntesis}
\begin{tcolorbox}[colback=piedra,colframe=ladrillo,boxrule=1pt,arc=3pt,fontupper=\small]
\textbf{\textcolor{cafe}{Israel en el Sinaí:}} un pueblo redimido antes de recibir la ley, llamado a ser \emph{segulá}, reino de sacerdotes y nación santa, con Dios habitando en medio de él. Rompe el pacto al pie del monte y en el desierto; Dios responde con su nombre misericordioso, preserva un remanente y anuncia destierro y regreso con un corazón circuncidado (Dt 4; 30). Desde ese horizonte escribe Pablo Romanos 9--11.
\end{tcolorbox}
\scriptsize\renewcommand{\arraystretch}{1.3}
\begin{tabular}{>{\RaggedRight\arraybackslash}p{3.0cm}>{\RaggedRight\arraybackslash}p{8.4cm}}
\toprule
{\color{ladrillo}\bfseries Eje} & {\color{ladrillo}\bfseries Lo que deja esta etapa}\\\midrule
Redención y vocación & Pascua, mar, Éx 19:4--6; Torá ante los pueblos (Dt 4:6)\\
Presencia y sacerdocio & Tabernáculo y Shejiná; Aarón, Leví, Israel entre las naciones\\
Restauración & Éx 34:6--7; Nm 14:20; Dt 4:25--31; 30:1--10; 32:36--43\\
Proyección & Isaías (nuevo éxodo, Espíritu), Jeremías (nuevo pacto), Oseas (Acor), Pablo (Ro 10--11)\\
\bottomrule
\end{tabular}
\elabpropia
\end{frame}

\begin{frame}{Hacia el Estudio 3}
\begin{columns}[T,onlytextwidth]
\begin{column}{0.48\textwidth}
\begin{cita}Se levantó después de ellos otra generación que no conocía a YHWH, ni la obra que él había hecho por Israel.\fuente{(Jue 2:10)}\end{cita}
\footnotesize La tierra fue dada (Jos 21:45), pero la generación siguiente olvidará el Éxodo y entrará en un ciclo.\par\medskip
\begin{tcolorbox}[colback=cafe,colframe=cafeosc,coltext=piedra,boxrule=0.8pt,arc=3pt,fontupper=\footnotesize]
\textbf{Estudio 3}\par\smallskip
De los Jueces a David: ciclos de apostasía y liberación, Rut la moabita, Samuel, Saúl y la unción de David.
\end{tcolorbox}
\end{column}
\begin{column}{0.48\textwidth}
\centering
\begin{tikzpicture}
\node[bloqueoro,minimum width=2.1cm] (a) at (90:1.55) {Apostasía};
\node[bloque,minimum width=2.1cm] (b) at (0:1.95) {Opresión};
\node[bloque,minimum width=2.1cm] (c) at (270:1.55) {Clamor};
\node[bloqueazul,minimum width=2.1cm] (d) at (180:1.95) {Liberación};
\draw[flecha] (a) to[bend left=25] (b);
\draw[flecha] (b) to[bend left=25] (c);
\draw[flecha] (c) to[bend left=25] (d);
\draw[flecha] (d) to[bend left=25] (a);
\node[font=\tiny\itshape,text=sepia,align=center] at (0,0) {el ciclo\\de Jueces};
\end{tikzpicture}
\elabpropia
\end{column}
\end{columns}
\end{frame}

% ===================================================
\section{Reflexión}
\begin{frame}{Para la reflexión}
\footnotesize
\begin{enumerate}\setlength{\itemsep}{6pt}
\item Dios dice «os he traído a mí» (Éx 19:4) antes de dar mandamiento alguno. ¿Qué cambia al leer la ley desde la redención?
\item Israel es \emph{segulá} de un Dios que dice «mía es toda la tierra». ¿Cómo evita esta frase que la elección se lea como desprecio de las naciones?
\item El nombre misericordioso (Éx 34:6--7) se revela tras el becerro. ¿Qué enseña sobre el pecado de Israel y la fidelidad de Dios en Romanos 11?
\item Pablo dice a gentiles que «nuestros padres» pasaron el mar (1~Co 10:1). ¿Qué responsabilidad trae recibir esa historia (Ro 11:18--21)?
\item ¿Qué dicen Rahab y los gabaonitas sobre el lugar de los gentiles en el pueblo de Dios?
\end{enumerate}
\end{frame}

% ===================================================
\section{Glosario}
\begin{frame}{Glosario}
\begin{columns}[T,onlytextwidth]
\begin{column}{0.48\textwidth}\scriptsize
\textbf{Émeq Ajor} (\heb{עֵמֶק עָכוֹר}): valle de la desgracia; puerta de esperanza (Os 2:15).\par\medskip
\textbf{Hilastérion} (\grc{ἱλαστήριον}): \emph{kapóret} en la LXX; Ro 3:25.\par\medskip
\textbf{Jérem} (\heb{חֵרֶם}): consagración al exterminio (Jos 6:17).\par\medskip
\textbf{Jésed} (\heb{חֶסֶד}): lealtad amorosa del pacto (Éx 34:6).\par\medskip
\textbf{Kapóret} (\heb{כַּפֹּרֶת}): tapa del arca (Éx 25:17).\par\medskip
\textbf{Mamléjet kohaním} (\heb{מַמְלֶכֶת כֹּהֲנִים}): reino de sacerdotes.\par\medskip
\textbf{Mishkán} (\heb{מִשְׁכָּן}): tabernáculo, morada.
\end{column}
\begin{column}{0.48\textwidth}\scriptsize
\textbf{Najalá} (\heb{נַחֲלָה}): heredad (Jos 13--21).\par\medskip
\textbf{Segulá} (\heb{סְגֻלָּה}): tesoro personal (Éx 19:5).\par\medskip
\textbf{Shavuot} (\heb{שָׁבוּעוֹת}): fiesta de las Semanas; Pentecostés.\par\medskip
\textbf{Shelosh esré midot} (\heb{שְׁלֹשׁ עֶשְׂרֵה מִדּוֹת}): trece atributos.\par\medskip
\textbf{Sukkot} (\heb{סֻכּוֹת}): fiesta de las Cabañas (Lv 23:43).\par\medskip
\textbf{Yam suf} (\heb{יַם־סוּף}): mar de los juncos.\par\medskip
\textbf{Yehoshúa} (\heb{יְהוֹשֻׁעַ}): Josué; \grc{Ἰησοῦς}.
\end{column}
\end{columns}
\end{frame}

% ===================================================
\section{Referencias}
\begin{frame}[allowframebreaks]{Referencias}
\tiny
\refapa{Barkay, G., Vaughn, A.~G., Lundberg, M.~J., \& Zuckerman, B. (2004). The amulets from Ketef Hinnom: A new edition and evaluation. \emph{Bulletin of the American Schools of Oriental Research, 334}, 41--71.}
\refapa{\emph{Biblia Hebraica Stuttgartensia} (5.ª ed.). (1997). Deutsche Bibelgesellschaft.}
\refapa{Bimson, J.~J. (1978). \emph{Redating the Exodus and conquest}. JSOT Press.}
\refapa{Bruins, H.~J., \& van der Plicht, J. (1995). Tell es-Sultan (Jericho): Radiocarbon results of short-lived cereal and multiyear charcoal samples from the end of the Middle Bronze Age. \emph{Radiocarbon, 37}(2), 213--220.}
\refapa{Callaway, J.~A. (1968). New evidence on the conquest of `Ai. \emph{Journal of Biblical Literature, 87}(3), 312--320.}
\refapa{Danby, H. (Trad.). (1933). \emph{The Mishnah}. Oxford University Press.}
\refapa{Epstein, I. (Ed.). (1935--1948). \emph{The Babylonian Talmud}. Soncino Press.}
\refapa{Faust, A. (2006). \emph{Israel's ethnogenesis: Settlement, interaction, expansion and resistance}. Equinox.}
\refapa{Finkelstein, I. (1988). \emph{The archaeology of the Israelite settlement}. Israel Exploration Society.}
\refapa{Gottwald, N.~K. (1979). \emph{The tribes of Yahweh}. Orbis Books.}
\refapa{Guggenheimer, H.~W. (1998). \emph{Seder Olam: The rabbinic view of biblical chronology}. Jason Aronson.}
\refapa{Hammer, R. (Trad.). (1986). \emph{Sifre: A Tannaitic commentary on the book of Deuteronomy}. Yale University Press.}
\refapa{Hays, R.~B. (1989). \emph{Echoes of Scripture in the letters of Paul}. Yale University Press.}
\refapa{Hoffmeier, J.~K. (2005). \emph{Ancient Israel in Sinai}. Oxford University Press.}
\refapa{Hoftijzer, J., \& van der Kooij, G. (1976). \emph{Aramaic texts from Deir `Alla}. Brill.}
\refapa{Josefo, F. (1987). \emph{The works of Josephus} (W. Whiston, Trad.). Hendrickson.}
\refapa{Kitchen, K.~A. (2003). \emph{On the reliability of the Old Testament}. Eerdmans.}
\refapa{Lauterbach, J.~Z. (Trad.). (1933--1935). \emph{Mekilta de-Rabbi Ishmael}. Jewish Publication Society.}
\refapa{Mendenhall, G.~E. (1954). Covenant forms in Israelite tradition. \emph{The Biblical Archaeologist, 17}(3), 50--76.}
\refapa{Mendenhall, G.~E. (1962). The Hebrew conquest of Palestine. \emph{The Biblical Archaeologist, 25}(3), 66--87.}
\refapa{Moo, D.~J. (1996). \emph{The epistle to the Romans}. Eerdmans.}
\refapa{Moran, W.~L. (Ed. y Trad.). (1992). \emph{The Amarna letters}. Johns Hopkins University Press.}
\refapa{\emph{Novum Testamentum Graece} (28.ª ed.). (2012). Deutsche Bibelgesellschaft.}
\refapa{ONE FOR ISRAEL Ministry. (2025, 3 de julio). \emph{What does «all Israel will be saved» mean in Romans 11?} \url{https://www.oneforisrael.org/what-does-all-israel-will-be-saved-mean-in-romans-11/}}
\refapa{Rahlfs, A., \& Hanhart, R. (Eds.). (2006). \emph{Septuaginta} (ed. rev.). Deutsche Bibelgesellschaft.}
\refapa{Sarna, N.~M. (1991). \emph{Exodus} (The JPS Torah Commentary). Jewish Publication Society.}
\refapa{Tov, E. (2012). \emph{Textual criticism of the Hebrew Bible} (3.ª ed.). Fortress Press.}
\refapa{Wood, B.~G. (1990). Did the Israelites conquer Jericho? A new look at the archaeological evidence. \emph{Biblical Archaeology Review, 16}(2), 44--58.}
\refapa{Younger, K.~L., Jr. (1990). \emph{Ancient conquest accounts}. JSOT Press.}
\refapa{Zuckerman, S. (2007). Anatomy of a destruction: Crisis architecture, termination rituals and the fall of Canaanite Hazor. \emph{Journal of Mediterranean Archaeology, 20}(1), 3--32.}
\end{frame}
\end{document}
